\documentclass{article}
\usepackage{pathfinder-paper}

\title{Priced Certificates for Conditional-CVaR Obedience}

\begin{document}
\maketitle

\begin{abstract}
For a fixed finite-signal policy, we study whether one realised batch of
conditional-loss samples can certify exact obedience under upper-tail
conditional value-at-risk (CVaR). A simultaneous confidence box and a
finite difference-constraint graph give either a supporting state-independent
reward, a negative-cycle witness, or abstention. Under a strict cycle margin,
the least nonnegative reward cap exceeds the true-utility oracle cap by at
most an explicit confidence term. Sampling cost can therefore be charged
alongside one reward payment. This is a local, confidence-qualified result;
it is not a competitive ratio for a learning-augmented system.
% ledger: 3, 8, 12, 14, 20
\end{abstract}

\section{Setting and provenance}

The supplied Q is a survey of learning-augmented algorithms. It proposes a
charge per predictor invocation and gives sufficient conditions for composing
component cost guarantees, including conditional guarantees at reachable
histories and relations among benchmarks \cite{zhao2026survey}. The supplied
P studies a one-type, finite-signal obedience test with state-independent
action rewards \cite{pathfinder2026reversal}. Its expected empirical CVaR
surrogate can reverse the true obedience decision at every fixed finite batch
size. P also proves deterministic cycle and common-reward tests when all
conditional utilities lie within a stated error bound. P explicitly leaves
one realised-batch certification outside those results. The mechanism paper
called Q inside P is distinct from the supplied Q.
% ledger: 4, 5, 20

We fix a finite signal set \(S\), a finite set \(B\) of available actions,
and a policy \(a:S\to B\) chosen before sampling. Every available action is
prescribed at some signal; alternatively, unprescribed actions are barred.
The agent has one type and evaluates an action after observing \(s\) by
\(U(b,s)=-\operatorname{CVaR}_{\beta}(L_{s,b})\), where
\(\beta\in(0,1]\) is upper-tail mass and \(L_{s,b}\in[0,H]\).
A state-independent reward \(r:B\to\mathbb R\) supports the policy when
\[
U(a(s),s)+r(a(s))\geq U(b,s)+r(b)
\quad\text{for every }s\in S\text{ and }b\in B.
\]
This reduced-form conditional-risk model excludes ex-ante CVaR of a plan,
report-stage incentives, and the system objective considered by Q.
% ledger: 4, 5, 8, 15, 20

The contribution below is the realised-data certificate and its priced,
local reward-cap comparison. P's graph theorem supplies the deterministic
step; the CVaR concentration step and margin reasoning have close published
antecedents discussed in Section~\ref{sec:related}.
% ledger: 7, 8, 12, 13, 14, 20

\section{A certificate from one realised batch}

For each of the \(N=|S||B|\) cells \((s,b)\), assume access to \(m\)
independent, identically distributed draws from its fixed conditional loss
law. Dependence between different cells is allowed. Write
\(\widehat U(b,s)\) for negative upper-tail CVaR of that cell's realised
empirical distribution. For \(0<\delta<1\), set
\[
e_m=\frac{H}{\beta}
\sqrt{\frac{\log(2N/\delta)}{2m}},\qquad
\mathcal E=\left\{
\max_{s,b}|\widehat U(b,s)-U(b,s)|\leq e_m
\right\}.
\]
The variational CVaR formula gives
\(\lvert\operatorname{CVaR}_{\beta}(F)
-\operatorname{CVaR}_{\beta}(G)\rvert
\leq H\|F-G\|_{\infty}/\beta\) for distributions on \([0,H]\).
The Dvoretzky--Kiefer--Wolfowitz inequality and a union bound yield
\(\Pr(\mathcal E)\geq1-\delta\). This bounds the error of a realised
batch relative to true CVaR; P's expected empirical CVaR is a different
quantity.
% ledger: 3, 8, 20

For distinct \(a,b\in B\), define the empirical robust edge and the true
edge by
\[
\begin{aligned}
w_m(a,b)&=\max_{s:a(s)=a}
\{\widehat U(b,s)-\widehat U(a,s)+2e_m\},\\
w^*(a,b)&=\max_{s:a(s)=a}
\{U(b,s)-U(a,s)\}.
\end{aligned}
\]
The graph has vertex set \(B\) and one such edge for each ordered pair of
distinct actions. A cycle's weight is the sum of its edge weights.
% ledger: 8, 12, 14, 20

\begin{theorem}[Realised-batch decision]\label{thm:certificate}
If the empirical robust graph has no positive cycle, its difference
constraints admit a vector \(r\) satisfying
\[
r(a)-r(b)\geq w_m(a,b)\quad(a\ne b).
\]
The recorded estimates, radius, graph and vector then certify exact true
obedience on \(\mathcal E\). If instead a displayed cycle of \(K\) edges
has empirical utility sum \(\widehat C\) satisfying
\(\widehat C+2Ke_m<0\), it certifies true infeasibility on
\(\mathcal E\). In every other case the procedure abstains.
% ledger: 3, 8, 14, 20
\end{theorem}

\begin{proof}
For every signal prescribing \(a\), the event \(\mathcal E\) implies
\(U(b,s)-U(a,s)\leq w_m(a,b)\). Summing difference constraints around
a cycle proves their necessity; P's path-potential construction proves
sufficiency when there is no positive cycle. The resulting \(r\)
satisfies every true obedience inequality on \(\mathcal E\).
Each utility difference in a cycle changes by at most \(2e_m\), so the
second inequality makes the corresponding true cycle negative. Such a cycle
cannot coexist with any supporting reward. The graph step is P's
deterministic rectangular-error test, applied to a random box with stated
coverage.
% ledger: 3, 8, 14, 20
\end{proof}

The certificate is sound with probability at least \(1-\delta\) over the
specified sample draw. A positive robust cycle alone proves neither true
infeasibility nor failure of a particular true reward: it can lead to
abstention. A negative-cycle test uses an upper confidence bound, not the
robust graph's failed feasibility check.
% ledger: 8, 20

\section{Margin, payment and sample price}

Assume every nontrivial simple true cycle in the graph has weight at most
\(-K\gamma\), where \(K\) is its length and \(\gamma>0\). This is
equivalent to a positive per-edge margin for the corresponding utility-cycle
test. On \(\mathcal E\), direct comparison of edge maxima gives
\[
w^*(a,b)\leq w_m(a,b)\leq w^*(a,b)+4e_m.
\]
Thus \(e_m<\gamma/4\) makes every empirical robust cycle negative, and
Theorem~\ref{thm:certificate} produces a reward witness. One sufficient
uniform batch size is
\[
m>\frac{8H^2}{\beta^2\gamma^2}
\log\frac{2N}{\delta}.
\]
A true cycle with utility sum at most \(-K\gamma\) likewise produces a
negative certificate under this radius. These are sufficient margin
conditions, not optimal sample-complexity claims.
% ledger: 8, 15, 17, 18, 20

We now require nonnegative rewards. Let \(n=|B|\), let \(R^*\) be the
smallest possible cap \(\max_a r(a)\) for a reward satisfying the true
constraints, and let \(R_m^*\) be the corresponding cap for the empirical
robust constraints. A cap is the maximum payment for one decision, not the
expected payment under a specified signal distribution.
% ledger: 10, 11, 12, 14, 18

\begin{proposition}[Local reward-cap comparison]\label{prop:cap}
Under the preceding cycle margin, on \(\mathcal E\) and when
\(e_m<\gamma/4\),
\[
R^*\leq R_m^*\leq R^*+4(n-1)e_m.
\]
If one conditional-loss draw costs \(\kappa_s\) in the decision
objective's units, sampling plus at most one reward payment is bounded by
\[
\kappa_sNm+R^*+4(n-1)e_m
\]
on the coverage event, apart from any separately charged graph computation.
% ledger: 10, 11, 12, 14, 18, 20
\end{proposition}

\begin{proof}
For any edge array without a positive cycle, the coordinatewise least
nonnegative solution of \(r(a)-r(b)\geq w(a,b)\) is the maximum weight
of a path starting at \(a\), including the empty path. A maximising path
can be simple, so it uses at most \(n-1\) edges. The edge comparison above
raises any such path weight by at most \(4(n-1)e_m\). Taking maximum
coordinates of these least solutions gives the stated cap comparison.
% ledger: 12, 14
\end{proof}

The payment term is necessary. For instance, edge requirements
\(r(a)-r(b)\geq M\) and
\(r(b)-r(a)\geq-M-\gamma\) have a negative cycle, yet any
nonnegative supporting reward has cap at least \(M\). A declared budget
\(R\) can be checked by adding \(0\leq r(a)\leq R\) to the robust
difference constraints. Failure of that budget check does not show that
the true array violates the budget.
% ledger: 10, 11, 12

For \(\kappa_s>0\), the sufficient upper bound in
Proposition~\ref{prop:cap} has variable part
\(\kappa_sNm+A m^{-1/2}\), where
\[
A=4(n-1)\frac{H}{\beta}
\sqrt{\frac{\log(2N/\delta)}{2}}.
\]
Its continuous unconstrained minimiser is
\(m_0=(A/(2\kappa_sN))^{2/3}\). The margin condition additionally
requires the strict batch-size inequality above; an integer batch is
chosen above that threshold and near the larger of the threshold and
\(m_0\). This only optimises the displayed sufficient bound.
% ledger: 17, 18

Q charges predictor invocations, whereas the bill here charges draws.
Treating one draw as one invocation, or one batch as one invocation, is an
extra accounting choice. Neither \(R^*\) nor
\(\kappa_sNm+R_m^*\) is Q's prediction-free offline system benchmark.
% ledger: 5, 8, 15, 20

\section{Two access boundaries}

The sample oracle must reveal losses for deviations at each signal, not
just losses for the prescribed action. In a two-signal swapped-action
instance, give each prescribed action constant loss \(3/2\) in both
worlds. Give every unplayed deviation loss \(2\) in one world and
\(1\) in another. Every on-policy transcript has the same law in both
worlds, for any sample size, while the true two-edge utility cycles have
weights \(+1\) and \(-1\). Thus on-policy logs alone cannot identify
the obedience decision. A generative conditional sampler, covered
exploration, or a justified identification model is a separate premise
whose cost must be counted.
% ledger: 9, 13, 15, 20

There is also no uniform finite-sample exact yes-or-no decision near the
cycle boundary. Let the safe loss be \(1\), risky loss be
\(2X\) with \(X\sim\operatorname{Bernoulli}(p)\), and
\(\beta=1/2\). In two signals, swap which named action is risky and
prescribe the risky action at each signal. For \(p<1/2\), risky CVaR is
\(4p\), giving true utility-cycle sum \(C(p)=2-8p\).
At \(p_{\pm}=1/4\pm\varepsilon\), with
\(0<\varepsilon<1/4\), the signs differ. Every fixed finite binary
transcript has positive probability in both worlds. A nonabstaining
sequential test with errors at most \(\delta<1/2\) in both worlds and
finite expected informative sample count \(T\) obeys
\[
\mathbb E_{p_-}T\,
\operatorname{kl}(p_-,p_+)
\geq\operatorname{kl}(1-\delta,\delta).
\]
The per-draw divergence is of order \(\varepsilon^2\), so expected
sampling cost grows as the margin shrinks. This obstruction permits
abstention and positive-margin promises.
% ledger: 6, 7, 15, 20

\section{Related work and interpretation}\label{sec:related}

Concentration for empirical CVaR is established by
Prashanth, Jagannathan and Kolla \cite{prashanth2019concentration}; a
Wasserstein treatment of empirical risk estimation that includes CVaR
appears in Prashanth and Bhat \cite{prashanth2019wasserstein}.
Agrawal, Koolen and Juneja \cite{agrawal2021tailrisk} give
fixed-confidence best-arm methods and confidence bounds for tail-risk
measures. These works publish the statistical main step used here. P
\cite{pathfinder2026reversal} supplies the deterministic graph step in
the supplied manuscript. Thus the contribution is their narrow
combination with a priced reward witness and local cap comparison, rather
than a new concentration or graph theorem.
% ledger: 7, 8, 13, 20

Chen \cite{chen2026persuasion} uses a positive risk-value margin to
retain strict incentive compatibility under approximation in a different
persuasion model. Balcan, Sandholm and Vitercik
\cite{balcan2019approximate} estimate approximate incentive compatibility
from samples for other mechanism classes. The sequential change-of-measure
ingredient behind the boundary example is standard in best-arm
identification \cite{kaufmann2016complexity}. None of these sources,
in the targeted searches recorded in \url{search.md}, was found to state
the specific conditional-CVaR reward-cap comparison above; that search
does not establish publication priority.
% ledger: 6, 7, 13, 20

\section{Limitations and open questions}

The local guarantee concerns a fixed policy, bounded conditional losses,
sample access for all compared cells, state-independent rewards and one
decision. It makes no claim for ex-ante CVaR, report-stage incentives,
unobserved counterfactuals, or a general learning-augmented pipeline.
The graph construction gives a high-confidence statement, not a pathwise
guarantee on failed coverage. A negative-cycle witness and a supporting
reward are sound only on the coverage event.
% ledger: 4, 5, 9, 15, 20

An expected-cost claim needs an abstention action and failure-path cost
control. If downstream cost \(D\leq B\) on the entire coverage event,
including abstention, and \(D\leq M\) on every path, then
\[
\mathbb E[\kappa_sNm+D]
\leq\kappa_sNm+B+\delta M.
\]
For Q's adaptive composition, both coverage and cost premises must hold
conditionally on each reachable pre-sampling history. A common positive
offline benchmark and component-to-system cost relations would also be
needed for a competitive ratio. None is supplied by this local result.
% ledger: 20, 22, 24

Open questions are how to price and justify counterfactual access, how to
allocate samples around critical graph cycles, and how transfers and a
bounded abstention fallback enter a common downstream objective. A first
system-level test would specify these elements in the two-signal setting
and prove, or refute, a conditional expected-cost comparison with its
prediction-free benchmark.
% ledger: 15, 20, 22, 24

\bibliographystyle{plainurl}
\bibliography{references}

\end{document}
